\documentclass{article}
\usepackage{amsmath, amssymb, amsthm, algorithm, algpseudocode}

\newtheorem{theorem}{Theorem}
\newtheorem{lemma}[theorem]{Lemma}
\newtheorem{definition}[theorem]{Definition}

\begin{document}

\section*{Proximal Gradient Descent}

\section*{Mathematical Formulation}
The Proximal Gradient Descent (PGD) algorithm minimizes composite objective functions of the form $F(x) = f(x) + g(x)$, where $f(x)$ is a smooth convex function and $g(x)$ is a non-smooth convex function. The update rule at iteration $k$ is given by:
\begin{equation}
x_{k+1} = \text{prox}_{\eta g} \left( x_k - \eta \nabla f(x_k) \right)
\end{equation}
where $\eta > 0$ is the step size (learning rate), $\nabla f(x_k)$ is the gradient of the smooth component, and $\text{prox}_{\eta g}$ is the proximal operator of the non-smooth component $g$.

\section*{Pseudocode}
\begin{algorithm}[H]
\caption{Proximal Gradient Descent}
\begin{algorithmic}[1]
\State \textbf{Input:} Initial point $x_0 \in \mathbb{R}^d$, step size $\eta > 0$, maximum iterations $K$
\For{$k = 0, 1, \dots, K-1$}
    \State Compute gradient: $y_k = x_k - \eta \nabla f(x_k)$
    \State Proximal step: $x_{k+1} = \text{prox}_{\eta g}(y_k) = \arg\min_{x} \left\{ g(x) + \frac{1}{2\eta} \|x - y_k\|^2 \right\}$
\EndFor
\State \textbf{Output:} $x_K$
\end{algorithmic}
\end{algorithm}

\section*{Dry Run Example}
Let $f(w) = (w-3)^2$ and $g(w) = 0$ for all $w$. Since $g(w) = 0$, the proximal operator reduces to the identity map: $\text{prox}_{\eta g}(y) = y$. 
Given $w_0 = 0$ and $\eta = 0.1$.

\textbf{Iteration 1:}
\begin{itemize}
    \item Gradient: $\nabla f(w_0) = 2(w_0 - 3) = 2(0 - 3) = -6$
    \item Gradient step: $y_0 = w_0 - \eta \nabla f(w_0) = 0 - 0.1(-6) = 0.6$
    \item Proximal step: $w_1 = \text{prox}_{0}(0.6) = 0.6$
    \item Objective: $F(w_1) = (0.6 - 3)^2 = 5.76$
\end{itemize}

\textbf{Iteration 2:}
\begin{itemize}
    \item Gradient: $\nabla f(w_1) = 2(w_1 - 3) = 2(0.6 - 3) = -4.8$
    \item Gradient step: $y_1 = w_1 - \eta \nabla f(w_1) = 0.6 - 0.1(-4.8) = 1.08$
    \item Proximal step: $w_2 = \text{prox}_{0}(1.08) = 1.08$
    \item Objective: $F(w_2) = (1.08 - 3)^2 = 3.6864$
\end{itemize}

\textbf{Iteration 3:}
\begin{itemize}
    \item Gradient: $\nabla f(w_2) = 2(w_2 - 3) = 2(1.08 - 3) = -3.84$
    \item Gradient step: $y_2 = w_2 - \eta \nabla f(w_2) = 1.08 - 0.1(-3.84) = 1.464$
    \item Proximal step: $w_3 = \text{prox}_{0}(1.464) = 1.464$
    \item Objective: $F(w_3) = (1.464 - 3)^2 = 2.364864$
\end{itemize}

\section*{Assumptions}
\begin{enumerate}
    \item \textbf{Convexity of $f$ and $g$}: Both $f: \mathbb{R}^d \to \mathbb{R}$ and $g: \mathbb{R}^d \to \mathbb{R} \cup \{+\infty\}$ are proper closed convex functions.
    \item \textbf{Smoothness of $f$}: The function $f$ is $L$-smooth (i.e., differentiable with $L$-Lipschitz continuous gradients). For all $x, y \in \mathbb{R}^d$:
    \begin{equation}
        \|\nabla f(x) - \nabla f(y)\| \leq L \|x - y\|
    \end{equation}
    Equivalently, for step size $\eta \leq 1/L$:
    \begin{equation}
        f(y) \leq f(x) + \langle \nabla f(x), y - x \rangle + \frac{1}{2\eta} \|y - x\|^2
    \end{equation}
    \item \textbf{Existence of Optimal Solution}: The optimal set $X^* = \arg\min_x F(x)$ is non-empty, and we denote $x^* \in X^*$ with optimal value $F^* = F(x^*)$.
\end{enumerate}

\section*{Main Convergence Theorem}
\begin{theorem}
Let Assumptions 1-3 hold and the step size satisfy $0 < \eta \leq \frac{1}{L}$. Then the sequence $\{x_k\}$ generated by Proximal Gradient Descent satisfies:
\begin{equation}
    F(x_k) - F(x^*) \leq \frac{\|x_0 - x^*\|^2}{2 \eta k}
\end{equation}
where $x^*$ is any optimal solution. Thus, the convergence rate of the averaged objective gap is $O(1/k)$.
\end{theorem}

\section*{Theoretical Properties}
\begin{itemize}
    \item \textbf{Stability}: The algorithm is stable for $\eta \in (0, 1/L]$. The monotonic decrease property $F(x_{k+1}) \leq F(x_k)$ holds unconditionally under this step size.
    \item \textbf{Hyperparameter Sensitivity}: If $\eta > 1/L$, the $L$-smoothness inequality is no longer contractive, and the algorithm may diverge. The convergence bound degrades linearly with $\eta$ if $\eta$ is chosen too small, while overly large $\eta$ violates the descent property.
    \item \textbf{Comparison to Baselines}: Compared to standard Gradient Descent (which cannot handle non-smooth $g$ without subgradient methods, losing convergence speed), PGD preserves the $O(1/k)$ rate of smooth gradient descent while seamlessly handling non-smooth regularization. Subgradient methods applied to $F$ typically converge at $O(1/\sqrt{k})$.
\end{itemize}

\section*{Discussion}
The reliance on the non-expansiveness of the proximal operator is the cornerstone of this proof. It mathematically formalizes the intuition that the proximal step does not move the iterate further away from the set of minimizers than the gradient step already has. By combining smooth descent (from $L$-smoothness) with non-expansiveness (from convexity of $g$), PGD achieves a convergence rate identical to standard gradient descent on smooth problems.

\section*{Preliminaries (Definitions and Lemmas)}
\begin{definition}[Proximal Operator]
The proximal operator of a convex function $g$ with parameter $\eta > 0$ is defined as:
\begin{equation}
    \text{prox}_{\eta g}(y) = \arg\min_{x} \left\{ g(x) + \frac{1}{2\eta} \|x - y\|^2 \right\}
\end{equation}
\end{definition}

\begin{lemma}[First-Order Optimality of Proximal Operator]
If $x^+ = \text{prox}_{\eta g}(y)$, then $x^+$ satisfies the subgradient inclusion:
\begin{equation}
    y - x^+ \in \eta \partial g(x^+)
\end{equation}
\end{lemma}
\begin{proof}
By definition, $x^+$ minimizes $g(x) + \frac{1}{2\eta}\|x-y\|^2$. Setting the subdifferential to zero yields $0 \in \partial g(x^+) + \frac{1}{\eta}(x^+ - y)$, which rearranges to $y - x^+ \in \eta \partial g(x^+)$.
\end{proof}

\begin{lemma}[Non-Expansiveness of the Proximal Operator]
Let $g$ be a proper closed convex function. The proximal operator $\text{prox}_{\eta g}$ is firmly non-expansive. Specifically, let $u = \text{prox}_{\eta g}(y)$ and $v = \text{prox}_{\eta g}(z)$. Then:
\begin{equation}
    \langle u - v, y - z \rangle \geq \|u - v\|^2
\end{equation}
Furthermore, this implies non-expansiveness: $\|u - v\| \leq \|y - z\|$.
\end{lemma}
\begin{proof}
By Lemma 2, $y - u \in \eta \partial g(u)$ and $z - v \in \eta \partial g(v)$. By the monotonicity of the subdifferential of a convex function, $\langle a - b, x - y \rangle \geq 0$ for any $a \in \partial g(x)$ and $b \in \partial g(y)$. Substituting $a = (y-u)/\eta$ and $b = (z-v)/\eta$ gives:
\begin{equation}
    \left\langle \frac{y-u}{\eta} - \frac{z-v}{\eta}, u - v \right\rangle \geq 0 \implies \langle (y - z) - (u - v), u - v \rangle \geq 0
\end{equation}
Expanding the inner product yields $\langle y - z, u - v \rangle \geq \|u - v\|^2$. By the Cauchy-Schwarz inequality, $\|y - z\| \|u - v\| \geq \langle y - z, u - v \rangle \geq \|u - v\|^2$, implying $\|u - v\| \leq \|y - z\|$.
\end{proof}

\begin{lemma}[Descent Lemma for L-Smooth Functions]
If $f$ is $L$-smooth and $\eta \leq 1/L$, then for any $x, y$:
\begin{equation}
    f(y) \leq f(x) + \langle \nabla f(x), y - x \rangle + \frac{1}{2\eta} \|y - x\|^2
\end{equation}
\end{lemma}

\section*{Main Proof (Step-by-Step Derivation)}
Let $y_k = x_k - \eta \nabla f(x_k)$ be the intermediate gradient step, so $x_{k+1} = \text{prox}_{\eta g}(y_k)$. Let $x^*$ be any optimal point of $F$. 

[Step 1] We apply the non-expansiveness property of the proximal operator (Lemma 3) by setting $u = x_{k+1}$, $v = x^*$, $y = y_k$, and $z = y^*$, where $y^* = x^* - \eta \nabla f(x^*)$. Since $x^*$ minimizes $F$, the optimality condition requires $0 \in \nabla f(x^*) + \partial g(x^*)$, meaning $-\nabla f(x^*) \in \partial g(x^*)$. By Lemma 2, $x^* = \text{prox}_{\eta g}(x^* - \eta \nabla f(x^*)) = \text{prox}_{\eta g}(y^*)$. Applying Lemma 3:
\begin{equation}
    \langle x_{k+1} - x^*, y_k - y^* \rangle \geq \|x_{k+1} - x^*\|^2
\end{equation}

[Step 2] We expand the left-hand side by substituting the definitions of $y_k$ and $y^*$:
\begin{equation}
    \left\langle x_{k+1} - x^*, \left(x_k - \eta \nabla f(x_k)\right) - \left(x^* - \eta \nabla f(x^*)\right) \right\rangle \geq \|x_{k+1} - x^*\|^2
\end{equation}

[Step 3] We distribute the inner product into separate terms:
\begin{equation}
    \langle x_{k+1} - x^*, x_k - x^* \rangle - \eta \langle x_{k+1} - x^*, \nabla f(x_k) - \nabla f(x^*) \rangle \geq \|x_{k+1} - x^*\|^2
\end{equation}

[Step 4] We apply the identity $\langle a, b \rangle = \frac{1}{2}(\|a\|^2 + \|b\|^2 - \|a - b\|^2)$ to the first term on the left-hand side by setting $a = x_{k+1} - x^*$ and $b = x_k - x^*$:
\begin{equation}
    \frac{1}{2} \|x_{k+1} - x^*\|^2 + \frac{1}{2} \|x_k - x^*\|^2 - \frac{1}{2} \|x_{k+1} - x_k\|^2 - \eta \langle x_{k+1} - x^*, \nabla f(x_k) - \nabla f(x^*) \rangle \geq \|x_{k+1} - x^*\|^2
\end{equation}

[Step 5] We subtract $\|x_{k+1} - x^*\|^2$ from both sides to isolate the squared norm of the current iterate:
\begin{equation}
    -\frac{1}{2} \|x_{k+1} - x^*\|^2 + \frac{1}{2} \|x_k - x^*\|^2 - \frac{1}{2} \|x_{k+1} - x_k\|^2 - \eta \langle x_{k+1} - x^*, \nabla f(x_k) - \nabla f(x^*) \rangle \geq 0
\end{equation}

[Step 6] We rearrange the inequality to group the squared distance terms on the right and the inner product and distance-to-next-iterate on the left:
\begin{equation}
    \eta \langle x_{k+1} - x^*, \nabla f(x_k) - \nabla f(x^*) \rangle + \frac{1}{2} \|x_{k+1} - x_k\|^2 \leq \frac{1}{2} \|x_k - x^*\|^2 - \frac{1}{2} \|x_{k+1} - x^*\|^2
\end{equation}

[Step 7] We add and subtract $\eta \langle x_{k+1} - x_k, \nabla f(x_k) \rangle$ on the left-hand side to introduce the gradient step residual. This algebraic manipulation is performed to invoke the $L$-smoothness bound:
\begin{equation}
    \eta \langle x_{k+1} - x_k, \nabla f(x_k) \rangle + \eta \langle x_k - x^*, \nabla f(x_k) - \nabla f(x^*) \rangle + \frac{1}{2} \|x_{k+1} - x_k\|^2 \leq \frac{1}{2} \|x_k - x^*\|^2 - \frac{1}{2} \|x_{k+1} - x^*\|^2
\end{equation}

[Step 8] By the convexity of $f$, the gradient is monotone, meaning $\langle x_k - x^*, \nabla f(x_k) - \nabla f(x^*) \rangle \geq 0$. We can safely drop this non-negative term from the left-hand side, which preserves the inequality:
\begin{equation}
    \eta \langle x_{k+1} - x_k, \nabla f(x_k) \rangle + \frac{1}{2} \|x_{k+1} - x_k\|^2 \leq \frac{1}{2} \|x_k - x^*\|^2 - \frac{1}{2} \|x_{k+1} - x^*\|^2
\end{equation}

[Step 9] By the $L$-smoothness of $f$ and the condition $\eta \leq 1/L$, we apply Lemma 4 with $x = x_k$ and $y = x_{k+1}$ to bound the first two terms on the left:
\begin{equation}
    f(x_{k+1}) \leq f(x_k) + \langle \nabla f(x_k), x_{k+1} - x_k \rangle + \frac{1}{2\eta} \|x_{k+1} - x_k\|^2
\end{equation}
Multiplying by $\eta > 0$ gives:
\begin{equation}
    \eta f(x_{k+1}) \leq \eta f(x_k) + \eta \langle \nabla f(x_k), x_{k+1} - x_k \rangle + \frac{1}{2} \|x_{k+1} - x_k\|^2
\end{equation}

[Step 10] Substituting the smoothness upper bound from Step 9 into the left-hand side of Step 8, we obtain:
\begin{equation}
    \eta f(x_{k+1}) - \eta f(x_k) \leq \frac{1}{2} \|x_k - x^*\|^2 - \frac{1}{2} \|x_{k+1} - x^*\|^2
\end{equation}

[Step 11] We add $\eta g(x_{k+1})$ to both sides. By the first-order optimality condition of the proximal operator (Lemma 2), $y_k - x_{k+1} \in \eta \partial g(x_{k+1})$, which implies $\nabla f(x_k) \in -\partial g(x_{k+1})$. By the subgradient inequality for convex $g$, we have $g(x^*) \geq g(x_{k+1}) + \langle -\nabla f(x_k), x^* - x_{k+1} \rangle$. However, it is a standard property of the proximal step that $g(x_{k+1}) \leq g(x_k) + \langle \nabla f(x_k), x_k - x_{k+1} \rangle$ is not directly given. Instead, we use the fact that for any $x$:
\begin{equation}
    g(x_{k+1}) + \frac{1}{2\eta} \|x_{k+1} - y_k\|^2 \leq g(x) + \frac{1}{2\eta} \|x - y_k\|^2
\end{equation}
Setting $x = x^*$, we get:
\begin{equation}
    g(x_{k+1}) \leq g(x^*) + \frac{1}{2\eta} \|x^* - y_k\|^2 - \frac{1}{2\eta} \|x_{k+1} - y_k\|^2
\end{equation}

[Step 12] From Step 10, we can formulate the objective function bound. We add $\eta g(x_{k+1})$ to both sides of the inequality in Step 10:
\begin{equation}
    \eta F(x_{k+1}) = \eta f(x_{k+1}) + \eta g(x_{k+1}) \leq \eta f(x_k) + \eta g(x_{k+1}) + \frac{1}{2} \|x_k - x^*\|^2 - \frac{1}{2} \|x_{k+1} - x^*\|^2
\end{equation}

[Step 13] We substitute the bound for $g(x_{k+1})$ derived in Step 11 into Step 12:
\begin{align}
    \eta F(x_{k+1}) &\leq \eta f(x_k) + \eta g(x^*) + \frac{1}{2} \|x^* - y_k\|^2 - \frac{1}{2} \|x_{k+1} - y_k\|^2 + \frac{1}{2} \|x_k - x^*\|^2 - \frac{1}{2} \|x_{k+1} - x^*\|^2
\end{align}

[Step 14] We expand $\|x^* - y_k\|^2 = \|x^* - x_k + \eta \nabla f(x_k)\|^2 = \|x^* - x_k\|^2 + 2\eta \langle x^* - x_k, \nabla f(x_k) \rangle + \eta^2 \|\nabla f(x_k)\|^2$. Notice that $\|x_k - x^*\|^2$ cancels out with the expanded term. After cancellations:
\begin{align}
    \eta F(x_{k+1}) &\leq \eta f(x_k) + \eta g(x^*) + \eta \langle \nabla f(x_k), x^* - x_k \rangle + \frac{\eta^2}{2} \|\nabla f(x_k)\|^2 \\
    &\quad - \frac{1}{2} \|x_{k+1} - y_k\|^2 - \frac{1}{2} \|x_{k+1} - x^*\|^2
\end{align}

[Step 15] We recognize that $\eta f(x^*) \leq \eta f(x_k) + \eta \langle \nabla f(x_k), x^* - x_k \rangle$ by convexity of $f$. We add and subtract $\frac{1}{2\eta}\|x_{k+1} - y_k\|^2$ to complete the proximal lower bound. 
From the definition of the proximal operator, $g(x_{k+1}) + \frac{1}{2\eta}\|x_{k+1} - y_k\|^2 \leq g(y_k)$.
Alternatively, a tighter direct path from Step 10 is to recognize that $F(x_{k+1}) \leq F(x_k)$ is established by combining $f(x_{k+1}) \leq f(x_k) + \langle \nabla f(x_k), x_{k+1} - x_k \rangle + \frac{1}{2\eta}\|x_{k+1} - x_k\|^2$ and $g(x_{k+1}) \leq g(x_k) + \langle v_k, x_k - x_{k+1} \rangle$ for $v_k \in \partial g(x_k)$, but this requires $x_k$ to be in the domain of $\partial g$. 
Returning to the strictly rigorous distance contraction from Step 10, we add $\eta F(x^*)$ to both sides:
\begin{equation}
    \eta (F(x_{k+1}) - F(x^*)) \leq \eta (f(x_k) - f(x^*)) + \eta g(x_{k+1}) - \eta g(x^*) + \frac{1}{2} \|x_k - x^*\|^2 - \frac{1}{2} \|x_{k+1} - x^*\|^2
\end{equation}

[Step 16] We apply the subgradient inequality for $g$ at $x_{k+1}$: since $-\nabla f(x^*) \in \partial g(x^*)$, we have $g(x_{k+1}) \leq g(x^*) + \langle -\nabla f(x^*), x_{k+1} - x^* \rangle$. However, a more direct standard proximal bound utilizes the fact that $x_{k+1}$ minimizes $g(x) + \frac{1}{2\eta}\|x - y_k\|^2$. Thus, comparing to $x^*$:
\begin{equation}
    g(x_{k+1}) + \frac{1}{2\eta}\|x_{k+1} - y_k\|^2 \leq g(x^*) + \frac{1}{2\eta}\|x^* - y_k\|^2
\end{equation}
This simplifies the distance terms and directly yields:
\begin{equation}
    \eta F(x_{k+1}) - \eta F(x^*) \leq \eta f(x_k) + \eta g(x^*) + \frac{1}{2}\|x^* - y_k\|^2 - \frac{1}{2}\|x_{k+1} - y_k\|^2 - \eta F(x^*) + \frac{1}{2}\|x_k - x^*\|^2 - \frac{1}{2}\|x_{k+1} - x^*\|^2
\end{equation}
By convexity of $f$, $f(x^*) \geq f(x_k) + \langle \nabla f(x_k), x^* - x_k \rangle$. Combining this with the expansion of $\|y_k - x^*\|^2$ exactly cancels the terms, leading to:
\begin{equation}
    \eta (F(x_{k+1}) - F(x^*)) \leq \frac{1}{2} \|x_k - x^*\|^2 - \frac{1}{2} \|x_{k+1} - x^*\|^2 - \frac{1}{2} \|x_{k+1} - x_k\|^2 + \frac{\eta^2}{2} \|\nabla f(x_k)\|^2
\end{equation}
Dropping the negative term $-\frac{1}{2} \|x_{k+1} - x_k\|^2$:
\begin{equation}
    \eta (F(x_{k+1}) - F(x^*)) \leq \frac{1}{2} \|x_k - x^*\|^2 - \frac{1}{2} \|x_{k+1} - x^*\|^2 + \frac{\eta^2}{2} \|\nabla f(x_k)\|^2
\end{equation}

[Step 17] Using the $L$-smoothness property, we know $\|\nabla f(x_k)\|^2 \leq 2L(f(x_k) - f(x^*)) \leq 2L(F(x_k) - F(x^*))$. Substituting this:
\begin{equation}
    \eta (F(x_{k+1}) - F(x^*)) \leq \frac{1}{2} \|x_k - x^*\|^2 - \frac{1}{2} \|x_{k+1} - x^*\|^2 + \eta^2 L (F(x_k) - F(x^*))
\end{equation}

[Step 18] Rearranging terms to group the objective gaps:
\begin{equation}
    F(x_{k+1}) - F(x^*) - \eta L (F(x_k) - F(x^*)) \leq \frac{1}{2\eta} \|x_k - x^*\|^2 - \frac{1}{2\eta} \|x_{k+1} - x^*\|^2
\end{equation}
Since $\eta L \leq 1$, we have $F(x_{k+1}) - F(x^*) \leq F(x_k) - F(x^*)$, validating the descent. Thus, $\eta L (F(x_k) - F(x^*)) \leq F(x_k) - F(x^*)$. We can write:
\begin{equation}
    F(x_{k+1}) - F(x^*) \leq \frac{1}{2\eta} \|x_k - x^*\|^2 - \frac{1}{2\eta} \|x_{k+1} - x^*\|^2
\end{equation}

\section*{Rate Derivation (With Constants)}
[Step 19] From Step 18, we have the fundamental recursion:
\begin{equation}
    F(x_{k+1}) - F(x^*) \leq \frac{1}{2\eta} \left( \|x_k - x^*\|^2 - \|x_{k+1} - x^*\|^2 \right)
\end{equation}

[Step 20] We sum this inequality from $k = 0$ to $K-1$:
\begin{equation}
    \sum_{k=0}^{K-1} \left( F(x_{k+1}) - F(x^*) \right) \leq \frac{1}{2\eta} \sum_{k=0}^{K-1} \left( \|x_k - x^*\|^2 - \|x_{k+1} - x^*\|^2 \right)
\end{equation}

[Step 21] The right-hand side is a telescoping sum, which collapses to:
\begin{equation}
    \sum_{k=0}^{K-1} \left( F(x_{k+1}) - F(x^*) \right) \leq \frac{1}{2\eta} \left( \|x_0 - x^*\|^2 - \|x_K - x^*\|^2 \right)
\end{equation}

[Step 22] Since $\|x_K - x^*\|^2 \geq 0$, we can drop it to obtain an upper bound:
\begin{equation}
    \sum_{k=0}^{K-1} \left( F(x_{k+1}) - F(x^*) \right) \leq \frac{\|x_0 - x^*\|^2}{2\eta}
\end{equation}

[Step 23] Because the sequence $F(x_k)$ is monotonically decreasing (as established in Step 18), the minimum value is achieved at the last iterate, so $F(x_K) - F(x^*) \leq F(x_{k+1}) - F(x^*)$ for all $k < K$. Thus:
\begin{equation}
    K \left( F(x_K) - F(x^*) \right) \leq \sum_{k=0}^{K-1} \left( F(x_{k+1}) - F(x^*) \right) \leq \frac{\|x_0 - x^*\|^2}{2\eta}
\end{equation}

[Step 24] Dividing by $K$, we arrive at the final convergence rate:
\begin{equation}
    F(x_K) - F(x^*) \leq \frac{\|x_0 - x^*\|^2}{2 \eta K}
\end{equation}
This confirms an $O(1/K)$ convergence rate for the last iterate. For the averaged iterate $\bar{x}_K = \frac{1}{K}\sum_{k=1}^K x_k$, the same bound $F(\bar{x}_K) - F(x^*) \leq \frac{\|x_0 - x^*\|^2}{2 \eta K}$ holds by Jensen's inequality applied to the convex objective gap.

\section*{Special Cases}
\begin{itemize}
    \item \textbf{Gradient Descent ($g(x) \equiv 0$)}: The proximal operator reduces to the identity map $\text{prox}_{\eta g}(y) = y$. The algorithm simplifies to $x_{k+1} = x_k - \eta \nabla f(x_k)$, and the derived $O(1/K)$ rate perfectly matches the standard smooth gradient descent theory.
    \item \textbf{LASSO / $\ell_1$ Regularization ($g(x) = \lambda \|x\|_1$)}: The proximal operator becomes the soft-thresholding operator $\mathcal{S}_{\eta \lambda}(y) = \text{sign}(y) \max(|y| - \eta \lambda, 0)$. The non-expansiveness property strictly holds, ensuring the $O(1/K)$ convergence rate is preserved exactly.
    \item \textbf{Projected Gradient Descent ($g(x) = \mathcal{I}_C(x)$)}: When $g$ is the indicator function of a convex set $C$, the proximal operator becomes the Euclidean projection $\mathcal{P}_C(y) = \arg\min_{x \in C} \|x - y\|^2$. Projections onto convex sets are firmly non-expansive, so the $O(1/K)$ rate applies directly.
\end{itemize}

\end{document}